\chapter[Python's dataset dictionary structure]{\python{'s} dataset dictionary structure}\label{sec:gretl4py:data:dict}
\python{'s'} dictionary, to be understood by \gretlpy,
must have the following entries:
\begin{itemize}
	\item \pythcode{varname}: A list containing the variable names.
	\item \pythcode{obs\_labels}: A list containing the observation labels.
	\item \pythcode{data\_array}: A \pkg{numpy} array (either column-major
	or row-major format) containing the observations for the variables.
	\item \pythcode{periodicity}: An integer indicating the data's periodicity,
	as defined in \cite{gretl_command}. A brief explanation is as follows:
	\begin{itemize}
		\item For time-series data:
		\begin{itemize}
			\item 1 = annual,
			\item 4 = quarterly,
			\item 12 = monthly,
			\item 52 = weekly,
			\item 5, 6, or 7 = daily,
			\item 24 = hourly.
		\end{itemize}
		\item For panel data: The periodicity represents the number
		of rows per data block. This corresponds to:
		\begin{itemize}
			\item The number of cross-sectional units in the case
			of stacked cross-sections, or
			\item The number of time periods in the case of
			stacked time-series data.
		\end{itemize}
		\item For simple cross-sectional data: The periodicity should
		be set to~1.
	\end{itemize}
\end{itemize}
